\documentclass[11pt,a4paper]{article}

\usepackage[margin=2.6cm]{geometry}
\usepackage{amsmath,amssymb}
\usepackage{booktabs}
\usepackage[colorlinks=true,linkcolor=black]{hyperref}

\title{\bf Cost of the marker scans:\\ memory and runtime versus system size}
\author{}
\date{\today}

\begin{document}
\maketitle

\begin{abstract}
\noindent
Measured scaling of the $Z_2$ marker calculation, used to set the \texttt{-{}-mem} and
\texttt{-{}-time} requests in the Slurm submit files. The Hilbert dimension is
$D = 4L^2$ and every point of a scan performs two dense Hermitian diagonalisations plus a
handful of dense $D\times D$ products, so \textbf{time grows as $L^6$ and peak memory as
$L^4$}. Both laws were fitted to end-to-end measurements at $L=8$--$28$ and reproduce them
to within a few percent.
\end{abstract}

\section{Where the cost comes from}

At each point of a scan the code builds the Hamiltonian, diagonalises it, forms the
projector $\rho$, builds $\widetilde{S}$, diagonalises that as well, flattens it, and
contracts it against the position operators for the marker. Counting only the
$\mathcal{O}(D^3)$ operations:

\begin{center}
\begin{tabular}{@{}llc@{}}
\toprule
step & operation & count \\
\midrule
\texttt{spectrum(H)}            & Hermitian eigendecomposition (with eigenvectors) & 1 \\
projector $\rho = UU^\dagger$   & dense product                                     & $\tfrac12$ \\
$\widetilde{S} = \rho S + S\rho - S$ & dense products                               & 2 \\
sign flattening                 & eigendecomposition + 2 products                   & 1 + 2 \\
marker                          & dense product                                     & 1 \\
\midrule
\multicolumn{2}{@{}l}{\emph{total}} & 2 eigh + $\sim$5.5 products \\
\bottomrule
\end{tabular}
\end{center}

\noindent
The two eigendecompositions dominate: measured on this machine a complex Hermitian
\texttt{eigh} costs about $0.8\,\mathrm{ns}\times D^3$ against $0.043\,\mathrm{ns}\times D^3$
for a \texttt{zgemm}, i.e.\ roughly $18\times$ more per floating point operation. Together
they account for $85$--$90\%$ of the runtime, which is why the optimisations available in the
matrix products would not change the picture much.

\paragraph{The $\theta$ scan.} When \texttt{optimize\_theta} is enabled, each candidate angle
adds one \emph{eigenvalue-only} solve, measured at $0.39$ of a full eigendecomposition. Since
$\widetilde{S}$ is linear in $S_\theta$, the two basis operators are built once per point and
reused, so $N_\theta$ candidates multiply the runtime by
\begin{equation}
  f_\theta \;=\; \frac{2 + 0.39\,N_\theta}{2} ,
  \label{eq:thetafactor}
\end{equation}
that is $1.39\times$ for $N_\theta=2$ and $1.58\times$ for $N_\theta=3$. Memory is unaffected:
the basis operators are freed before the flattening allocates.

\section{Measurements}

End-to-end, one point per process, on 11 cores (Apple Silicon, arm64):

\begin{center}
\begin{tabular}{@{}rrrr@{}}
\toprule
$L$ & $D = 4L^2$ & peak RSS [GB] & time per point [s] \\
\midrule
 8 &  256 & 0.134 &  0.09 \\
12 &  576 & 0.196 &  0.37 \\
16 & 1024 & 0.324 &  1.21 \\
20 & 1600 & 0.609 &  4.20 \\
24 & 2304 & 1.134 & 10.60 \\
28 & 3136 & 1.983 & 24.45 \\
\bottomrule
\end{tabular}
\end{center}

\section{Fitted laws}

\paragraph{Memory.} About twelve $D\times D$ complex128 arrays are live at the peak, inside
the sign flattening, on top of a constant interpreter and kwant footprint:
\begin{equation}
  \boxed{\;\text{peak}\,[\mathrm{GB}] \;\simeq\; 0.12 \;+\; 3072\times10^{-9}\,L^4\;}
  \label{eq:mem}
\end{equation}
This reproduces the measurements to better than $2\%$ across the fitted range: at $L=20$ it
predicts $0.612$ GB against $0.609$ measured, and at $L=28$ it predicts $2.01$ GB against
$1.98$ measured.

\paragraph{Time.} Anchoring the $L^6$ law at the largest measured size,
\begin{equation}
  \boxed{\;t_{\rm point} \;\simeq\; 24.45\ \mathrm{s} \times \left(\frac{L}{28}\right)^{6}
  \times f_\theta\;}
  \label{eq:time}
\end{equation}
This is mildly \emph{optimistic} below $L\approx20$, where the matrices are too small to
saturate the BLAS and the measured times exceed the law (at $L=20$: $4.2$ s measured against
$3.25$ s predicted). Since those sizes run in minutes either way, the discrepancy is
irrelevant in practice; the law is accurate where it matters.

\section{Extrapolation}

Per point, and per scan of $N$ points, from \eqref{eq:mem} and \eqref{eq:time}:

\begin{center}
\begin{tabular}{@{}rrrrrr@{}}
\toprule
$L$ & $D$ & peak RSS & $t_{\rm point}$ & $W$ scan ($N=30$, $f_\theta=1$)
& $M$ scan ($N=50$, $f_\theta=1.58$) \\
\midrule
10 &   400 &  0.15 GB & 0.05 s & 2 s     & 4 s \\
20 &  1600 &  0.61 GB & 3.3 s  & 1.6 min & 4.3 min \\
30 &  3600 &  2.61 GB & 36 s   & 18 min  & 48 min \\
40 &  6400 &  7.98 GB & 203 s  & 1.7 h   & 4.5 h \\
50 & 10000 & 19.32 GB & 775 s  & 6.5 h   & 17.0 h \\
75 & 22500 & 97.32 GB & 2.5 h  & 74 h    & 193 h \\
\bottomrule
\end{tabular}
\end{center}

\noindent
These are times for \emph{one} disorder realisation. A batch of $N_{\rm samples}$ realisations
is that many independent jobs, not a longer one.

\section{Slurm requests}

The submit generators apply a $1.5\times$ margin to memory and $3\times$ to walltime, then
round up to a sensible step. For the $M$ scan with $N_\theta = 3$:

\begin{center}
\begin{tabular}{@{}rrll@{}}
\toprule
$L$ & predicted peak & \texttt{-{}-mem} & \texttt{-{}-time} \\
\midrule
10 &  0.15 GB &  2G & 0-00:30:00 \\
20 &  0.61 GB &  2G & 0-00:30:00 \\
30 &  2.61 GB &  4G & 0-04:00:00 \\
40 &  7.98 GB & 12G & 1-00:00:00 \\
50 & 19.32 GB & 32G & 3-00:00:00 \\
\bottomrule
\end{tabular}
\end{center}

\noindent
Two points on the Slurm side. The request is \texttt{-{}-mem} (total for the task) rather than
\texttt{-{}-mem-per-cpu}, because the threads share one set of matrices instead of each
holding a copy; with \texttt{-{}-cpus-per-task=8} a setting of \texttt{-{}-mem-per-cpu=1G}
would silently cap the job at 8 GB, which $L=40$ and $L=50$ would exceed. And the thread count
must be pinned to the allocation (\texttt{OMP\_NUM\_THREADS} and friends), otherwise the BLAS
spawns one thread per core it can see on the node rather than per core granted, which
oversubscribes badly.

\section{Caveats}

\begin{itemize}
  \item The absolute constants are for 11 Apple Silicon cores. Cluster nodes will differ;
        the $3\times$ walltime margin is meant to absorb that, but a short run should be
        checked before trusting the large sizes.
  \item \texttt{eigh} parallelises decently but far from linearly, so doubling the cores
        does not halve the time. When the queue is busy, more jobs with fewer cores each
        will usually clear faster than fewer wide ones.
  \item $L=75$ is included only to show where the wall is: at $\sim\!97$ GB and days per
        single realisation it needs a large-memory node, and disorder averaging on top is
        not realistic without moving away from full diagonalisation.
\end{itemize}

\end{document}
